\chapter{Architettura del software}
\label{cap:architettura}

\section{Moduli}

Il codice è organizzato in moduli distinti in base alle rispettive responsabilità. Ogni sottocartella di src/ raccoglie le funzionalità relative a uno specifico aspetto dell'applicazione.
\begin{table}[H]
\centering
\footnotesize
\caption{Struttura di \file{src/}.}
\label{tab:moduli}
\begin{tabular}{@{}p{0.20\linewidth}p{0.72\linewidth}@{}}
\toprule
\textbf{Modulo} & \textbf{Responsabilit\`a} \\
\midrule
\file{common/} & definizione del tipo scalare (\code{scalar\_t}), configurabile come \code{float} o \code{double}, tolleranza di
  validazione, utility di terminazione e cronometro \\
\file{index/} & partizionamento a blocchi degli indici globali: dimensione,
  offset e proprietario di ogni blocco \\
\file{gen/} & generatore riproducibile senza stato (\cref{sec:generazione}) \\
\file{serial/} & implementazione seriale minimale, usata come oracolo \\
\file{kernel/} & \code{local\_gemm}: un file per implementazione, \code{.c} o
  \code{.cu}, tutte dietro la stessa interfaccia \\
\file{mpi/} & griglia cartesiana, layout dei blocchi, distribuzione dei dati,
  algoritmo distribuito $Y = AX$ \\
\file{bench/} & driver di misura e validazione, parsing delle opzioni, output CSV \\
\bottomrule
\end{tabular}
\end{table}

Il tipo \code{scalar\_t} permette di utilizzare la stessa base di codice sia in singola precisione (\code{float}, FP32) sia in doppia precisione (\code{double}, FP64); la precisione viene selezionata a tempo di compilazione.

Il modulo \file{mpi/} non dipende direttamente da uno specifico backend
di calcolo: per l'esecuzione del prodotto locale utilizza esclusivamente
l'interfaccia comune definita in \file{kernel/kernel.h}. Di conseguenza, l'algoritmo
MPI rimane invariato sia utilizzando un'implementazione CPU sia utilizzando uno dei backend CUDA.
La scelta del backend avviene quindi in fase di compilazione,
senza richiedere modifiche al codice che gestisce la distribuzione dei dati e le comunicazioni MPI.